% GENERATED FILE. Edit canonical lessons, then run npm run generate:books.
% canonical-source-hash: fnv1a64:e6477246fd430193
% canonical-lessons: ES-C736-frente, ES-C736-mejilla, ES-C736-barbilla, ES-C736-barba, ES-C736-oreja

\chapter{La cara}
\label{ch:es-la-cara}
\hlchaptermodality{736}

\begin{chapteropening}
\textbf{By the end of this chapter:} \emph{I can describe a face: the forehead, cheeks, chin, beard and ears.}

Complete the last lesson of chapter 736: I can describe a face: the forehead, cheeks, chin, beard and ears.
\end{chapteropening}

\section[la frente]{la frente — the forehead}

\label{lesson:ES-C736-frente}



\begin{quote}
Before the new one: say the Spanish for the brain, then the Spanish for to beat.
\end{quote}

\subsection*{You'll want to know: la frente}
\textbf{la frente} — \textquotedblleft{}the forehead\textquotedblright{}.

\begin{quote}
Tienes la frente caliente.
\end{quote}

\begin{quote}
\emph{Your forehead is hot.}
\end{quote}

\emph{La frente} is the forehead; \emph{el frente} is the front.

\begin{cousinweb}
\emph{frente} is inherited from Latin \emph{frōns}, \emph{frontem} \textquotedblleft{}forehead, front\textquotedblright{} (EtymDB), through Old Spanish \emph{fruente}.
\end{cousinweb}

\begin{grammarlens}[title={What you've built}]
One more word for telling, explaining and describing at B1.
\end{grammarlens}

\subsection*{Your turn}
\begin{itemize}
\raggedright
  \item \emph{Say it:} \emph{la frente}
  \item \emph{Say it:} \emph{la frente}, once more
  \item \emph{Say it:} a sentence of your own with \emph{la frente}
  \item \emph{Recall:} say the Spanish for the brain, then the Spanish for to beat, then say \emph{la frente} again
\end{itemize}

\begin{tcolorbox}[breakable,colback=teal!4,colframe=teal!35!black,title={Before you move on}]
What does la frente mean? (\textquotedblleft{}The forehead\textquotedblright{}.) Say it once more.
\end{tcolorbox}
\section[la mejilla]{la mejilla — the cheek}

\label{lesson:ES-C736-mejilla}



\begin{quote}
Before the new one: say the Spanish for to beat, then the Spanish for the forehead.
\end{quote}

\subsection*{You'll want to know: la mejilla}
\textbf{la mejilla} — \textquotedblleft{}the cheek\textquotedblright{}.

\begin{quote}
Me dio un beso en la mejilla.
\end{quote}

\begin{quote}
\emph{She kissed me on the cheek.}
\end{quote}

Colloquially also \emph{el cachete} in Latin America.

\begin{cousinweb}
\emph{mejilla} is usually traced to Latin \emph{maxilla} \textquotedblleft{}jaw\textquotedblright{}; the Latin \emph{-x-} became the Spanish \emph{j} sound.
\end{cousinweb}

\begin{grammarlens}[title={What you've built}]
One more word for telling, explaining and describing at B1.
\end{grammarlens}

\subsection*{Your turn}
\begin{itemize}
\raggedright
  \item \emph{Say it:} \emph{la mejilla}
  \item \emph{Say it:} \emph{la mejilla}, once more
  \item \emph{Say it:} a sentence of your own with \emph{la mejilla}
  \item \emph{Recall:} say the Spanish for to beat, then the Spanish for the forehead, then say \emph{la mejilla} again
\end{itemize}

\begin{tcolorbox}[breakable,colback=teal!4,colframe=teal!35!black,title={Before you move on}]
What does la mejilla mean? (\textquotedblleft{}The cheek\textquotedblright{}.) Say it once more.
\end{tcolorbox}
\section[la barbilla]{la barbilla — the chin}

\label{lesson:ES-C736-barbilla}



\begin{quote}
Before the new one: say the Spanish for the forehead, then the Spanish for the cheek.
\end{quote}

\subsection*{You'll want to know: la barbilla}
\textbf{la barbilla} — \textquotedblleft{}the chin\textquotedblright{}.

\begin{quote}
Tiene una cicatriz en la barbilla.
\end{quote}

\begin{quote}
\emph{He has a scar on his chin.}
\end{quote}

\emph{Levantar la barbilla} — to lift your chin. \emph{El mentón} means the same.

\begin{cousinweb}
\emph{barbilla} is built on \emph{barba} \textquotedblleft{}beard, chin\textquotedblright{} with the small ending \emph{-illa}; \emph{barba} is inherited from Latin \emph{barba} (EtymDB).
\end{cousinweb}

\begin{grammarlens}[title={What you've built}]
One more word for telling, explaining and describing at B1.
\end{grammarlens}

\subsection*{Your turn}
\begin{itemize}
\raggedright
  \item \emph{Say it:} \emph{la barbilla}
  \item \emph{Say it:} \emph{la barbilla}, once more
  \item \emph{Say it:} a sentence of your own with \emph{la barbilla}
  \item \emph{Recall:} say the Spanish for the forehead, then the Spanish for the cheek, then say \emph{la barbilla} again
\end{itemize}

\begin{tcolorbox}[breakable,colback=teal!4,colframe=teal!35!black,title={Before you move on}]
What does la barbilla mean? (\textquotedblleft{}The chin\textquotedblright{}.) Say it once more.
\end{tcolorbox}
\section[la barba]{la barba — the beard; the chin}

\label{lesson:ES-C736-barba}



\begin{quote}
Before the new one: say the Spanish for the cheek, then the Spanish for the chin.
\end{quote}

\subsection*{You'll want to know: la barba}
\textbf{la barba} — \textquotedblleft{}the beard; the chin\textquotedblright{}.

\begin{quote}
Mi padre lleva barba.
\end{quote}

\begin{quote}
\emph{My father has a beard.}
\end{quote}

\emph{Llevar barba} — to have a beard.

\begin{cousinweb}
\emph{barba} is inherited from Latin \emph{barba} \textquotedblleft{}beard\textquotedblright{} (EtymDB); English \emph{barber} comes from the same word.
\end{cousinweb}

\begin{grammarlens}[title={What you've built}]
One more word for telling, explaining and describing at B1.
\end{grammarlens}

\subsection*{Your turn}
\begin{itemize}
\raggedright
  \item \emph{Say it:} \emph{la barba}
  \item \emph{Say it:} \emph{la barba}, once more
  \item \emph{Say it:} a sentence of your own with \emph{la barba}
  \item \emph{Recall:} say the Spanish for the cheek, then the Spanish for the chin, then say \emph{la barba} again
\end{itemize}

\begin{tcolorbox}[breakable,colback=teal!4,colframe=teal!35!black,title={Before you move on}]
What does la barba mean? (\textquotedblleft{}The beard; the chin\textquotedblright{}.) Say it once more.
\end{tcolorbox}
\section[la oreja]{la oreja — the ear (outer)}

\label{lesson:ES-C736-oreja}



\begin{quote}
Before the new one: say the Spanish for the chin, then the Spanish for the beard.
\end{quote}

\subsection*{You'll want to know: la oreja}
\textbf{la oreja} — \textquotedblleft{}the ear (outer)\textquotedblright{}.

\begin{quote}
Lleva un pendiente en la oreja.
\end{quote}

\begin{quote}
\emph{He wears an earring in his ear.}
\end{quote}

\emph{La oreja} is the outer ear; \emph{el oído} is hearing and the inner ear.

\begin{cousinweb}
\emph{oreja} is inherited from Latin \emph{auricula} \textquotedblleft{}little ear\textquotedblright{} through spoken \emph{oricla} (EtymDB), a small form of \emph{auris} \textquotedblleft{}ear\textquotedblright{}.
\end{cousinweb}

\begin{grammarlens}[title={What you've built}]
One more word for telling, explaining and describing at B1.
\end{grammarlens}

\subsection*{Your turn}
\begin{itemize}
\raggedright
  \item \emph{Say it:} \emph{la oreja}
  \item \emph{Say it:} \emph{la oreja}, once more
  \item \emph{Say it:} a sentence of your own with \emph{la oreja}
  \item \emph{Recall:} say the Spanish for the chin, then the Spanish for the beard, then say \emph{la oreja} again
\end{itemize}

\begin{tcolorbox}[breakable,colback=teal!4,colframe=teal!35!black,title={Before you move on}]
What does la oreja mean? (\textquotedblleft{}The ear (outer)\textquotedblright{}.) Say it once more.
\end{tcolorbox}
